% ===================================================================
%  CADRO N.º 13 — O hotel. — O restaurante. O café.
%  Traduction 2026 d'apres texto/io, controlee sur texto/fr, texto/es
%  et texto/pt. Consigne : voir texto/gl/10-cadro-01.tex.
%
%  DEUX NOTES D'AUTEUR, l'une et l'autre en gras-italique dans le
%  fac-simile, comme dans l'ido : la premiere renvoie au tableau 6
%  pour le detail de la chambre, la seconde au vocabulaire pour la
%  carte des plats.
%
%  LE TABLEAU DE L'HOTEL PARTAGE SON LEXIQUE EN DEUX A L'INTERIEUR
%  D'UNE MEME PIECE, et c'est le troisieme cas apres la montagne du
%  tableau 9 et la mer du tableau 10. Ce qui SE MANGE est galicien de
%  bout en bout — « manteiga » le beurre, « queixo » le fromage,
%  « peixe » le poisson, « camaróns » les crevettes, « froitas » les
%  fruits, « espinacas », « callos », « perna » le gigot — et ce qui
%  SE SERT arrive avec le service : « menú », « restaurante »,
%  « ascensor », « telegrama », « telefonar », « billar »,
%  « automóbil », « triciclo ».
%
%  ET LE MOT LE PLUS PARLANT DU TABLEAU EST « PROPINA ». Le francais
%  dit « pourboire », l'ido calque « gratifiketo » sur le latin de la
%  gratification ; le galicien, l'espagnol et le portugais disent
%  tous trois « propina », qui est le latin de ce qu'on boit EN PLUS.
%  Trois langues qui gardent ensemble un mot que ni l'ecole ni le
%  chemin de fer n'ont apporte : il vient de l'auberge, qui est plus
%  vieille que les deux.
% ===================================================================

%%K t13-noto-1 noto
\VUnotes{143.2mm}{%
\textit{\VUgras{(*) Sobre os detalles do cuarto, véxase o cadro N.\textsuperscript{o}
6.}}%
}

%%K t13-apar-1 apar
\VUsaut{3.0mm}
\VUtitre{15.0pt}{60}{CADRO N.\textsuperscript{o} 13}
\VUsaut{1.6mm}
\VUpk{10.2pt}{40}{O hotel. --- O restaurante.}
\VUsaut{2.0mm}
\VUpk{10.2pt}{40}{O café.}
\VUsaut{3.4mm}

%%K t13-01-1 p
1. --- Chegado onte pola noite a Royan, aloxeime no «\textit{Hotel do Océano}», que me
recomendara un amigo.

%%K t13-01-2 p
Déronme un belo \VUgras{cuarto} \textsuperscript{(2)} no segundo \VUgras{andar}
\textsuperscript{(1)}, onde durmín moi ben (*).

%%K t13-02-1 p
2. --- Esta mañá, ao saír do meu cuarto, no que a \VUgras{criada}
\textsuperscript{(3)} me trouxera o almorzo, entrei no \VUgras{fumadoiro}
\textsuperscript{(5)}, que se usa tamén como \VUgras{sala de lectura}
\textsuperscript{(4)}, para ler os \VUgras{xornais} \textsuperscript{(6)} fumando un
\VUgras{cigarro} \textsuperscript{(8)}. Despois de ler, baixei polo
\VUgras{ascensor} \textsuperscript{(9)} e fun pasear pola vila.

%%K t13-03-1 p
3. --- Como esquecera preguntarlle ao \VUgras{empregado} \textsuperscript{(12)} do
hotel a que hora se serven as comidas na \VUgras{mesa común} \textsuperscript{(11)},
volvín cedo e aproveitei para ir bañarme no \VUgras{cuarto de baño}
\textsuperscript{(10)} do hotel. Despois fun outra vez á \VUgras{caixa}
\textsuperscript{(15)} para telefonar a un dos meus amigos. Pero o \VUgras{teléfono}
\textsuperscript{(13)} estaba ocupado; ao ver o meu apuro, a \VUgras{caixeira}
\textsuperscript{(14)}, que estaba a rexistrar unha \VUgras{factura}
\textsuperscript{(17)} no \VUgras{libro de viaxeiros} \textsuperscript{(16)},
pediulle ao \VUgras{porteiro} \textsuperscript{(18)} que mandase o \VUgras{botones}
\textsuperscript{(19)} facer o meu recado \VUgras{en bicicleta}
\textsuperscript{(20)}.

%%K t13-04-1 p
4. --- Deille as grazas e saín para darlle unha propina ao \VUgras{mozo}
\textsuperscript{(24)} (o traballador de forza) que trouxera da estación no seu
\VUgras{carro de man} \textsuperscript{(25)} a miña \VUgras{caixa de sombreiros}
\textsuperscript{(26)}, \VUgras{a miña maleta} \textsuperscript{(27)} e \VUgras{o meu
saco de viaxe} \textsuperscript{(28)}. O \VUgras{carteiro} \textsuperscript{(21)},
que acababa de chegar, deume unha \VUgras{carta} \textsuperscript{(22)}.

%%K t13-05-1 p
5. --- Quedei aínda un pouco na soleira da porta, preto do \VUgras{caixa do correo}
\textsuperscript{(23)}, mirando o movemento da rúa. O \VUgras{condutor}
\textsuperscript{(30)} do \VUgras{ómnibus} \textsuperscript{(29)} cargaba no seu coche
o \VUgras{baúl} \textsuperscript{(31)} dun \VUgras{viaxeiro} \textsuperscript{(32)} ao
que acompañaba o \VUgras{hoteleiro} \textsuperscript{(33)}. Un novo
\VUgras{telegrafista} \textsuperscript{(35)} traía sen présa un \VUgras{telegrama}
\textsuperscript{(34)} para un cliente do hotel; un \VUgras{turista}
\textsuperscript{(36)} coa súa \VUgras{máquina de fotografía} \textsuperscript{(38)}
en bandoleira consultaba a súa \VUgras{guía} \textsuperscript{(37)}.

%%K t13-tit-1 sub
\VUsaut{4.0mm}
\VUpk{10.2pt}{40}{Hora do xantar.}
\VUsaut{3.0mm}

%%K t13-06-1 p
6. --- Diante da reixa, un \VUgras{recadeiro} \textsuperscript{(39)} despreocupado
dormitaba entre o seu \VUgras{gancho de carga} \textsuperscript{(40)} e a súa
\VUgras{caixa de betún} \textsuperscript{(41)}, mentres unha nova
\VUgras{lavandeira} \textsuperscript{(42)} que levaba un \VUgras{cesto}
\textsuperscript{(43)} de roupa ría do xesto solemne do \VUgras{maître}
\textsuperscript{(44)} que estaba de pé xunto á porta.

%%K t13-07-1 p
7. --- Ao ver o \VUgras{cociñeiro} \textsuperscript{(45)}, que saía da súa
\VUgras{cociña} \textsuperscript{(47)} situada no \VUgras{soto} \textsuperscript{(48)}
para mandar un \VUgras{pinche} \textsuperscript{(46)} facer un recado, lembrei que a
hora do xantar estaba preto e entrei no restaurante, atravesando o patio no que un
\VUgras{chofer} \textsuperscript{(50)} gardaba o seu \VUgras{automóbil}
\textsuperscript{(49)}, mentres un \VUgras{mecánico} \textsuperscript{(52)} amañaba,
segundo as indicacións do \VUgras{ciclista} \textsuperscript{(53)}, un
\VUgras{triciclo de petróleo} \textsuperscript{(51)} que acababa de ter un accidente.

%%K t13-08-1 p
8. --- Pregunteille a un novo \VUgras{botones} \textsuperscript{(54)} que levaba uns
\VUgras{vasos de cervexa} \textsuperscript{(55)} se a mesa común estaba baixo a
galería e, como me respondeu que si, sentei nunha das mesas e fixen que me servisen
unha comida excelente.

%%K t13-noto-2 noto
\VUnotes{143.2mm}{%
\textit{\VUgras{(*) Coa axuda da lista de pratos posta no vocabulario, o mestre
poderá\nl variar doadamente o menú de abaixo.}}%
}

%%K t13-tit-2 sub
\VUsaut{4.0mm}
\VUpk{10.2pt}{40}{Un xantar excelente.}
\VUsaut{3.0mm}

%%K t13-09-1 p
9. --- O camareiro tróuxome o menú (*) do xantar e a carta dos viños. Escollín e
pedín como \VUgras{entremés camaróns} con \VUgras{manteiga,} e, como
\VUgras{entrada, callos á moda de Caen.} Non comín \VUgras{peixe,} pero pedín unha
porción de \VUgras{perna} de año e, como \VUgras{prato de legumes, espinacas} con
nata. Despois, como ningún dos \VUgras{pratos intermedios} me gustou, fixen traer
sobremesas variadas, \VUgras{queixo,} \VUgras{froitas} e \VUgras{biscoitos.} Engadín
ademais un excelente \VUgras{viño} de Saint-Émilion e, moi contento coa miña comida,
deixei a mesa despois de darlle ao \VUgras{camareiro} \textsuperscript{(57)} unha bela
\VUgras{propina} \textsuperscript{(59)}.

%%K t13-10-1 p
10. --- Ao verme disposto a marchar, o \VUgras{xerente} \textsuperscript{(60)}
propúxome ir ao balcón para admirar o espléndido panorama do \VUgras{Océano}
\textsuperscript{(62)}. Ao lonxe albiscábanse \VUgras{barquiñas} de pesca
\textsuperscript{(63)}, \VUgras{veleiros} \textsuperscript{(64)} e un enorme
\VUgras{paquebote} \textsuperscript{(65)} cuxa silueta negra se recortaba sobre as
\VUgras{nubes} brancas \textsuperscript{(67)} do \VUgras{horizonte}
\textsuperscript{(66)}. Unha brisa lixeira facía ondear o \VUgras{estandarte}
\textsuperscript{(70)} chantado por riba do \VUgras{letreiro} \textsuperscript{(69)}
do \VUgras{vestíbulo} \textsuperscript{(68)}.

%%K t13-tit-3 sub
\VUsaut{3.0mm}
\VUpk{10.2pt}{40}{Divertimentos.}
\VUsaut{3.0mm}

%%K t13-11-1 p
11. --- O espectáculo era tan interesante que decidín subir mañá á terraza do café
levando comigo uns \VUgras{binóculos} \textsuperscript{(71)} ou un
\VUgras{catalexo} \textsuperscript{(72)}.

%%K t13-12-1 p
12. --- Ao deixar o balcón, fun ao café do hotel para tomar unha \VUgras{bebida}
\textsuperscript{(73)} xogando unha \VUgras{partida de billar} \textsuperscript{(75)}.
Atravesei sen parar a sala do café, na que moitos consumidores xogaban ao
\VUgras{xadrez} \textsuperscript{(78)}, ás \VUgras{damas} \textsuperscript{(79)}, ao
\VUgras{dominó} \textsuperscript{(80)}, ao \VUgras{chaquete} \textsuperscript{(81)}
ou ás \VUgras{cartas} \textsuperscript{(82)}, e subín directamente á \VUgras{sala de
billar} \textsuperscript{(74)}.

%%K t13-13-1 p
13. --- Cando rematou a partida, baixei ao andar baixo e pedinlle ao camareiro un
\VUgras{sobre} \textsuperscript{(84)}, \VUgras{papel} \textsuperscript{(83)}, un
\VUgras{selo de correos} \textsuperscript{(85)}, unha \VUgras{pluma}
\textsuperscript{(87)} e \VUgras{tinta} \textsuperscript{(86)} para escribir unha
carta.

%%K t13-13-2 p
Despois chamei un \VUgras{cocheiro} \textsuperscript{(92)} cuxo \VUgras{coche}
\textsuperscript{(88)} parara diante do café. Pregunteille polo prezo por horas ou por
carreira e, cando nos puxemos de acordo no prezo, abriume a \VUgras{porta do coche}
\textsuperscript{(89)} e instaloume. Subiu de novo ao \VUgras{pescante}
\textsuperscript{(91)}, botou a andar os cabalos axiña co \VUgras{látego}
\textsuperscript{(93)} e marchamos para un paseo encantador.
